\subsection{Vector Fields}

\content{
\begin{definition} Let $D$ be a set in $\R^2$. A vector field on $D$ is a
function $\myvec{F}$ which assigns to each point in $D$, a two dimensional
vector $\myvec{F}(x,y)$.
\end{definition}
}{% presentation
\vspace{2in}
}{% nopresentation
}{% comments
Draw a sketch of how we visualize these. 
}

\content{
\begin{definition} Let $E$ be a subset of $\R^3$. A vector field on $E$ is a
function $\myvec{F}$ which to each point of $E$ assigns a three dimensional
vector $\myvec{F}(x,y,z)$.
\end{definition}
}{% presentation
\vspace{2in}
}{% nopresentation
}{% comments
sketch again. 
}

\content{
\begin{example} Sketch the vector field on $\R^2$ defined by 
$\myvec{F}(x,y) = \la -y, x\ra. $. 
\end{example}
}{% presentation
\vspace{2in}
}{% nopresentation
}{% comments
Mention that this is frequently denoted 
$$\myvec{F}(x,y) = -y\myvec{i}+x\myvec{j}. $$
}

\content{
\begin{example} Sketch the vector field given by $\myvec{F}(x,y,z) =
z\myvec{k}$. 
\end{example}
}{% presentation
\vspace{2in}
}{% nopresentation
}{% comments
}


\subsubsection{Gradient Fields}
\content{
Recall that if $f$ is a function of two variables $x$ and $y$, then the gradient
of $f$ is defined to be
$$ \nabla f(x,y) = \la f_x(x,y), f_y(x,y)\ra. $$
Thus, $\nabla f$ is really a vector field, and is called a gradient vector
field. 
}{% presentation
}{% nopresentation
}{% comments
}

\content{
\begin{example} Sketch the gradient vector field of the function $f(x,y) =
x^2y-y^3. $
\end{example}
}{% presentation
}{% nopresentation
}{% comments
}
